\section{Introducción}

Este documento ha sido estructurado siguiendo rigurosamente las directrices del Manual de Publicación de la American Psychological Association en su séptima edición \parencite{apa2020}. La arquitectura del proyecto ha sido modularizada para facilitar la redacción científica colaborativa y la intervención organizada de agentes de inteligencia artificial.

\subsection{Planteamiento del Problema}
En la elaboración de manuscritos académicos de alta exigencia, la integración de contenido, formato tipográfico y gestión bibliográfica suele concentrarse en archivos únicos monolíticos. Como destaca \textcite{knuth1984tex}, los sistemas de composición tipográfica como \TeX{} fueron concebidos para independizar la presentación formal del contenido intelectual. La adopción de una arquitectura modular permite que distintos roles de trabajo o agentes especializados operen sobre segmentos delimitados sin riesgo de corrupción del código maestro.

\subsection{Objetivos de la Investigación}
El propósito central de este marco de trabajo es proporcionar una base de redacción estandarizada que permita:
\begin{enumerate}
    \item Garantizar el cumplimiento del 100\% de las normas APA 7 (márgenes, tipografía, citas y jerarquía de encabezados).
    \item Facilitar la coordinación y trazabilidad entre agentes de IA para las fases de redacción, citación y auditoría de estilo.
    \item Simplificar la compilación mediante \texttt{biblatex} y el motor \texttt{biber}.
\end{enumerate}
